\section{Introduction}
\label{sec:intro}

This paper studies solution counting for ordered Combinatorial Discretizable Distance Geometry Problem (Combinatorial DDGP) instances with pruning edges. Fix a predecessor system and let \(B=V\setminus V_0\) be the binary Branch-and-Prune decision set, where \(V_0=\{v_1,\ldots,v_K\}\) is the initial clique. Pruning edges determine a constrained branch set \(B_{\mathrm{c}}\subseteq B\); the remaining \(f=\lvert B\setminus B_{\mathrm{c}}\rvert\) branch decisions are free. We prove that, under a generic mirror-separation hypothesis and relative to any reference solution, the feasible codes on \(B_{\mathrm{c}}\) form an affine space over \(\mathbb{F}_2\). Consequently, the number of feasible full branch codes is
\[
2^{
f+
\operatorname{rank}_{\mathbb{F}_2}\begin{bmatrix} M \\ V \end{bmatrix}
-
\operatorname{rank}_{\mathbb{F}_2}(V)
},
\]
where \(M\) is the graph-derived branch-generator matrix and \(V\) is the labeled violation matrix for pruning constraints.

The Distance Geometry Problem (DGP) asks for a geometric representation of a weighted graph in $\mathbb{R}^K$: given a simple undirected graph $G = (V, E)$ and an edge weight function $d: E \to \mathbb{R}_+$, find an embedding $x: V \to \mathbb{R}^K$ such that $\|x_i - x_j\|_2 = d_{ij}$ for every $\{i,j\}\in E$~\cite{mucherino2012discretizable}. The problem appears in protein structure determination from nuclear magnetic resonance (NMR) distance data~\cite{liberti2008branch}, sensor network localization~\cite{abud2018k}, and the geometry of nanostructures~\cite{mucherino2020analysis}.

When a vertex order on $V = \{1, \ldots, n\}$ is part of the input, the search space can be discretized. In DDGP and DMDGP instances~\cite{mucherino2012discretizable,lavor2012discretizable}, each non-seed vertex $j > K$ has a predecessor set $U_j \subset \{1, \ldots, j-1\}$ of size $K$. Once the predecessors are placed, $K$-lateration determines the position of $j$ as the intersection of $K$ spheres centered at the predecessor coordinates~\cite{liberti2014number}. Generically, this intersection contains at most two points, organizing realizations as paths in the Branch-and-Prune (BP) tree~\cite{liberti2008branch}.

We naturally partition the edges in $G$ into two sets:
\begin{enumerate}
    \item \textbf{Discretization edges} ($E_D = \{\{i, j\} \in E \mid i \in U_j\}$), which we use to perform the $K$-lateration steps and branch into the binary tree.
    \item \textbf{Pruning edges} ($E_P = E \setminus E_D$), which represent additional distance constraints that the BP algorithm uses to check the feasibility of branched positions and prune inconsistent paths in the tree.
\end{enumerate}

General DDGP counting cannot be purely graph-combinatorial in the unrestricted sense: numerical edge weights can change the number of feasible branches~\cite{abud2024impossible}. In the Combinatorial DDGP, the predecessor cliques restore enough reflection structure to make a rank count possible after pruning constraints are labeled by their relevant mirrors. The generic mirror-separation assumption excludes the proper algebraic exceptional set where a geometrically violating transformation accidentally preserves a distance. Outside this set, the feasible branch differences are exactly the zero-violation shifts.

Throughout the paper, all branch codes, vector spaces, sums, kernels, and ranks are over the finite field $\mathbb{F}_2$.

Our contribution is threefold:
\begin{enumerate}
    \item We decompose the branch decisions into pruning-constrained and pruning-free bits, and represent constrained branch transformations using graph-derived cone generators and base generators.
    \item We introduce a labeled violation matrix that separates violations caused by different mirror cliques, keeping track of mirror-specific incompatibilities.
    \item We prove the \emph{Rank-Count Theorem}: relative to any reference solution $s^\ast$, the feasible constrained branch codes form an affine coset with direction $\mathcal{K}_F$, the space of zero-violation shifts. Consequently,
    \[
    \lvert \Xi \rvert
    =
    2^{f+\operatorname{rank}_{\mathbb{F}_2}\begin{bmatrix} M \\ V \end{bmatrix} - \operatorname{rank}_{\mathbb{F}_2}(V)},
    \]
    where $M$ is the generator mask matrix and $V$ is the labeled violation matrix.
\end{enumerate}

Binary branch-code representations are standard for paths in the Branch-and-Prune tree; related binary representations have been widely used to analyze symmetry and degrees of freedom in distance geometry~\cite{mucherino2012exploiting,mucherino2020analysis}.

We organize the remainder of the paper as follows. Section~\ref{sec:related} reviews DMDGP symmetries, Combinatorial DDGP counts, DDGP recognition, and the counting barrier for general DDGP instances. Sections~\ref{sec:subproblems}--\ref{sec:shifts} define branch codes, generators, labeled violations, and feasible branch shifts. Section~\ref{sec:results} proves the main theorem. Section~\ref{sec:example} gives a fully worked 6-vertex example. Section~\ref{sec:scope} discusses scope, genericity, and the no-pruning Combinatorial DDGP special case.
